\section{Agents Across Scientific Domains}
\label{sec:domains}

This section compares representative systems through their research objects, workflow roles, decision mechanisms, interaction and feedback, human responsibilities, and scientific outcomes. Tables~\ref{tab:representatives-a}--\ref{tab:representatives-c} provide a compact reference; the analyses below explain differences within each domain. Object-specific evidence determines whether a work appears in more than one domain, following Section~\ref{sec:methods}.

\subsection{Formal, Information, and Computational Sciences}

Formal and computational investigations provide executable artifacts and relatively explicit evaluation procedures. FunSearch uses a language model to propose programs, executes them against a researcher-specified evaluator, and retains selected programs to guide subsequent generation~\cite{asd0512}. Its cap-set application concerns a concrete combinatorial object. Researchers supply the evaluator, initial program, and program skeleton, while the system evolves a critical function. The resulting scientific evidence concerns the constructions obtained through these programs and depends on the evaluator's correspondence to the mathematical question.

Research-workflow systems delegate a broader set of activities. The AI Scientist connects idea development, code-based experiments, manuscript production, and automated reviewing in machine-learning settings~\cite{asd0069}. Agent Laboratory organizes literature review, experimentation, and reporting around a human-provided research idea, with feedback permitted throughout the workflow~\cite{asd0018}. These systems extend language-based coordination across retrieval, coding, and assessment, whereas FunSearch concentrates on proposal generation coupled to external evaluation and selection.

The search representation provides a more specific basis for comparing program-search systems. FunSearch usually evolves a critical function within a supplied skeleton; AlphaEvolve permits changes across multiple functions and can co-evolve a candidate construction with the search procedure used to improve it~\cite{asd0512,asd0063}. Both retain researcher-provided evaluation, but the latter exposes more of the algorithm to revision. AlphaEvolve's internal ablations examine evolutionary selection, context, model choice, and, for matrix multiplication, full-file evolution under common computational-budget axes~\cite{asd0063}. These tests inform attribution within the evaluated tasks. Differences in tasks, models, and resources prevent a direct ranking of the two systems from their headline results.

Mathematical constructions also reveal how an output becomes new knowledge. In the eleven-dimensional kissing-number problem, AlphaEvolve produced 593 nonzero integer-coordinate points satisfying specified norm and pairwise-distance conditions. A geometric lemma maps this finite witness to a valid kissing configuration, improving the previous lower bound of 592~\cite{asd0063}. The result is supported by the checkable construction and its mathematical implication. It establishes a stronger lower bound, rather than the exact kissing number, and publication of verification code is distinct from completed independent reproduction.

\subsection{Physics, Astronomy, and Cosmic Sciences}

Physical-science agents connect research objectives to calibration, measurement, and inference. In the self-driving quantum laboratory k-agents, language-based procedure interpretation is coupled to state-machine execution on superconducting quantum hardware~\cite{asd0090}. Measurement reports determine transitions between experimental states, including calibration retries and explicit failure outcomes. The apparatus and supplied experimental procedures constrain the admissible actions; evaluation concerns physical outcomes as well as successful execution.

Astronomical systems place decisions elsewhere in the workflow. StarWhisper Telescope links observation planning, telescope operation, image processing, transient detection, and follow-up observation proposals~\cite{asd0136}. A multi-agent cosmological-analysis system instead coordinates retrieval, coding, execution, and parameter inference from observations; an administrator or human checks code before execution~\cite{asd0121}. These examples distinguish acquisition of time-sensitive observations from inference under specified likelihood and parameter assumptions.

A recent preprint extends the comparison from procedural control to explanatory-model revision. In the AI Theorist study of $\alpha$-RuCl$_3$, human-acquired reflection and photocurrent spectra supplied the observations to be explained. The agents compared computational routes, rejected a tested model that could not jointly explain the gap and optical features, and revised a constrained electronic state that predicted an unobserved dominant peak. Changes to orbital occupations, screening, and electron--hole treatment yielded candidate excitonic assignments and contrasting polarization predictions~\cite{asd1081}. The delegated step was selection of a new physical calculation in response to a discrepancy. The polarization dependence remained a prediction for future measurement; human theoretical review supplied additional assessment without establishing independent experimental confirmation.

Transfer between instruments reveals a further distinction between a reusable strategy and its local realization. An agent-assisted Grover-state preparation strategy was developed on one three-qubit NMR processor and applied sequentially to four others sharing a common interface~\cite{asd1064}. Each deployment rebuilt local Hamiltonian models, measurement references, and control pulses. Spectral overlap required a revised fitting procedure, subsequently reused with new measurements on another processor. Operators prepared and adjusted instruments and authorized changes. This preprint demonstrates adaptation within one processor family, with substantial local reconstruction, rather than portability of fixed settings across quantum platforms.

\subsection{Chemical Sciences}

Chemical agents connect written procedures, experimental selection, and analytical measurements. The synthesis-literature digitization system translates descriptions into machine-executable protocols, permits expert correction, and combines corrected procedures with simulated execution before synthesis~\cite{asd0600}. Its contribution is the implementation of specified chemistry through reproducible operations, rather than an inference of new chemical behavior from protocol completion.

Reactivity-seeking systems delegate the choice of subsequent experiments. The organic-synthesis robot uses learned predictions and spectral changes to explore reaction space, while a chemist manually follows up predictions associated with the reported reaction discoveries~\cite{asd0734}. Reactify combines robotic execution and analytical feedback with a neural reactivity-selection model, learning from expert-labelled spectra and retraining as further reactions are observed~\cite{asd0428}. Here, feedback changes experimental priorities, rather than only checking a prescribed procedure.

Exploratory selection can also change the phenomena available for investigation. Grizou and colleagues' formulation robot selected mixtures to reach temporary targets in a researcher-defined space of droplet behaviours, rather than maximize one fixed property. At 27~$^\circ$C, curiosity-driven exploration covered $73.4\pm15.2\%$ of the measured behaviour space, compared with $22.5\pm2.1\%$ for random formulation sampling under the same 1,000-experiment budget~\cite{asd0575}. An unexpected difference between repeated campaigns prompted the researchers to investigate temperature, which had initially been neither controlled nor recorded. Subsequent temperature-controlled experiments, longer trajectories, NMR measurements, and chemical interventions connected the newly observed motion patterns to dissolution and hydrolysis processes. The agent exposed a wider range of phenomena; identifying the omitted variable and developing the explanation required researcher-led follow-up. Together with manual isolation and characterization after reactivity screening~\cite{asd0734}, this case shows that discovery depends on how a selected signal becomes a separately investigated chemical phenomenon or identity.

\subsection{Materials Science}

Materials agents investigate composition--structure--property relationships through computational planning and measured optimization. MAPPS combines workflow planning, tool-code generation, and a scientific mediator for crystal generation, structure prediction, and property-constrained candidate generation; scientist feedback enters through the mediator~\cite{asd0001}. Its computational environment supports adaptation of the investigation, while candidate assessment remains separate from synthesis and physical characterization.

Experimental optimizers obtain feedback from material measurements. The mobile robotic chemist uses Bayesian search within a researcher-defined photocatalyst-mixture space, selecting subsequent experiments from measured hydrogen-production outcomes~\cite{burger2020}. Otto couples Bayesian optimization to aqueous-electrolyte measurements, including electrochemical stability~\cite{asd0487}. Researchers define ingredients, feasible operations, and objectives before the automated selection process.

A closer comparison separates the optimized property from its downstream use. Otto searched aqueous salt mixtures for electrochemical stability, whereas Clio optimized ionic conductivity in a non-aqueous electrolyte space~\cite{asd0487,asd0417}. Clio's 42-experiment campaign was followed by human electrolyte transfer and pouch-cell testing of six selected formulations against an a priori baseline. This supplies a tested connection from a screening objective to fast-charging performance under the reported cell conditions. The paper's approximately sixfold search-acceleration estimate, however, came from simulated campaigns over a fitted objective, rather than a matched physical random-search campaign~\cite{asd0417}. The two studies address related formulation decisions but different chemical spaces and endpoints; their optimum values cannot rank the platforms.

Optimization data can also become useful knowledge after the original optimizer is removed. In semi-autonomous sputter growth of $\beta$-Ga$_2$O$_3$, Bayesian optimization minimized optically derived Urbach energy, an indicator of absorption-tail disorder. Researcher-led random-forest analysis summarized the sampled response as largely additive with a residual temperature--oxygen interaction, motivating sequential one-variable tuning and local two-variable refinement. A human experimenter used this rule without an ML optimizer and reached 163~meV, compared with the initial optimization's 182~meV~\cite{asd0999}. Low-disorder conditions from both routes also supported growth on another substrate configuration within the tested window. The contribution is an experimentally tested local process rule, not a complete microscopic growth mechanism. Its usefulness outside the original acquisition path supplies stronger evidence than fit to the optimization data alone.

\subsection{Earth and Environmental Sciences}

Observation acquisition is a central agent role in Earth and environmental research. The Autonomous Earth-Observing Sensorweb connects event detection to observation requests, planning, and follow-up acquisition~\cite{asd0722}. Science teams specify interests and priorities before detections trigger requests and replanning. The agent mechanism allocates sensing resources as new information becomes available; its scientific contribution depends on the environmental observations acquired.

OASIS evaluates geological targets against scientist-specified preferences and resource constraints, connecting target priority to rover planning and follow-up actions. Its integrated opportunistic-science loop was demonstrated with the FIDO rover in the JPL Mars Yard~\cite{asd0691}. OpenEarth-Agent instead combines planning, coding, checking, and tool creation in benchmarked Earth-observation data-processing workflows~\cite{asd0650}. These cases distinguish model-based acquisition planning from language-based coordination of available data and computational tools.

Operational acquisition, ground-tested reactive observation, and benchmarked analysis establish different outcomes. Improved target selection or task execution requires separate evidence before supporting an improved interpretation of environmental processes. Environmental and mission-operation objects can both justify domain assignments when separately evaluated; satellite data use alone does not establish mission autonomy.

\subsection{Life Sciences}

Biological agents participate in hypothesis testing, perturbation selection, and data analysis. Robot Scientist uses a symbolic biological model to connect functional-genomic hypotheses with experiment design and yeast growth assays~\cite{king2004}. Measured growth discriminates among model predictions, while experimental preparation and plate handling still involve people. The 2004 study reconstructed known gene functions, demonstrating integrated hypothesis testing rather than discovery of an unknown function.

BioDiscoveryAgent combines biological knowledge, previous outcomes, and tools to choose genetic-perturbation experiments~\cite{asd0020}. Its evaluation retrieves outcomes from previously collected perturbation datasets, rather than reporting a new prospective experimental campaign. CellAgent uses hierarchical agent decisions, execution checks, and iterative evaluation to select analysis tools and settings for supplied single-cell data~\cite{asd0138}. These language-based systems address experimental design and analytical coordination under different access to biological evidence.

Comparison on a common selection task gives a more direct account of the contribution. PerTurboAgent re-evaluated an adapted BioDiscoveryAgent and active-learning baselines on eleven predefined K562 gene-expression phenotypes. Across ten retrospective rounds, its mean cumulative hit ratio was 0.440, compared with 0.240 for the adapted BioDiscoveryAgent and 0.255 for random selection augmented by enrichment analysis~\cite{asd0034}. A hit was a perturbation producing a sufficiently strong expression change in the specified gene programme; the denominator was the known set of hits in the existing screen. These values therefore concern recovery of phenotype-specific hits, not prospective disease benefit. Researchers supplied curated data and phenotype definitions. Removing enrichment and gene-set enrichment analysis reduced the ratio to 0.371, whereas removing the predictive ML action yielded 0.433~\cite{asd0034}. The ablations implicate expression analysis and prior knowledge as important components. Different tool and action configurations prevent assigning the whole advantage to planning alone, and the study's re-run comparison must be distinguished from scores in BioDiscoveryAgent's original, differently defined experiments.

\subsection{Medical and Health Sciences}

Medical agents support tasks ranging from electronic-health-record analysis to therapeutic-candidate design. HealthFlow combines planning, execution, evaluation, and reflective memory for healthcare research~\cite{asd0002}. Evaluation differs across tasks: ensemble LLM scoring is used for EHRFlowBench, human task-completion assessment for MedAgentBoard, and answer correctness for medical-reasoning benchmarks. These endpoints concern analytical and research-task performance, with the evaluator forming part of the conditions under which the result holds.

Tippy organizes agents around the design--make--test--analyze cycle and demonstrates platform integration through constructed use cases, including analytical feedback such as HPLC retention time~\cite{asd0008}. Virtual Lab combines a human researcher, specialist agents, and scientific tools to design SARS-CoV-2 nanobody candidates that undergo experimental assessment~\cite{asd0054}. Human researchers guide the collaboration and human biologists perform validation.

The transition from candidate generation to biological evidence is explicit in Virtual Lab. Starting from four characterized nanobodies, its computational pipeline selected 92 mutant candidates; subsequent expression and binding experiments identified two with newly acquired or improved binding profiles beyond the starting sequence backgrounds~\cite{asd0054}. The resulting contribution is a set of experimentally assessed molecules and binding phenotypes, without establishing a new therapeutic mechanism or clinical efficacy.

Reusing an analytical workflow provides a different kind of extension. SAGE's main multimodal biomarker example analysed FABP5 expression and tertiary-lymphoid-structure abundance in TCGA bladder-cancer data; an institutional-cohort example evaluated a different LAG3 hypothesis~\cite{asd0828}. The latter therefore does not constitute external replication of the former biomarker. Distinct datasets and hypotheses demonstrate workflow reuse, while transport of the same measurement relation or clinical association requires a corresponding comparison. This distinction is essential when interpreting the biological and clinical reach of an agent's outputs.

\subsection{Agricultural, Food, Forestry, Animal, and Fishery Sciences}

The current agricultural subset includes food hypotheses, plant phenotyping, and disease analysis. FoodPuzzle evaluates food--molecule relationship predictions against existing database evidence~\cite{asd0695}. PhenoAssistant supports conversational multi-agent phenotype extraction, visualization, and analysis, but its tool-selection evaluation measures the reasonableness of tool chains rather than the correctness of their outputs~\cite{asd0510}. These cases distinguish grounded hypothesis assessment from the validation of an analytical result.

Aleks supports iterative data-driven research in a plant-disease setting. Its authors report that manual verification detected code conflating training and test performance, and that another run contained transient data leakage~\cite{asd0634}. RAVA implements a different mechanism: a recurrent policy trained with proximal policy optimization selects informative visual observations for plant-disease diagnosis~\cite{asd1133}. Its image-dataset evaluation supports sequential visual inference, rather than prospective intervention in a field environment.

RAVA's explicit cross-dataset tests also delimit the robustness claim. With training on PlantVillage and testing on PlantDoc under the paper's most severe corruption protocol, RAVA achieved 18.78\% accuracy and ResNet-50 11.74\%. When training included both datasets, the corresponding results were 6.57\% and 11.27\%, respectively~\cite{asd1133}. The same table thus reports low absolute performance and a change in ordering between training protocols. These image-level tests support a conditional comparison of visual policies; they do not establish camera movement, field intervention, or agronomic benefit. Data distribution and training conditions determine whether sequential attention improves the particular diagnostic task.

\subsection{Engineering and Industrial Technology Sciences}

Engineering agents translate design objectives into physical tests or numerical investigations. BEAR, a Bayesian experimental autonomous researcher, connects fabrication, testing, model updates, and subsequent design selection using measured structural toughness~\cite{asd0504}. The tested structure family and objective are supplied by researchers, defining the scope of the measured optimization result.

Foam-Agent uses LLM-based coordination for mesh preparation, input generation, execution, debugging, and visualization in computational fluid dynamics~\cite{asd0047}. It reports execution success over 110 cases, with selected examples additionally comparing generated configurations and velocity fields against expert-prepared simulations. MechAgents coordinates problem formulation, coding, execution, and correction in mechanics~\cite{asd0128}. In one elasticity example, code executed successfully but extracted von Mises stress instead of the requested shear-stress component; a subsequent human instruction corrected the error.

A follow-up BEAR study provides a comparison within the same structural-design problem. Finite-element predictions supplied a prior, while physical experiments learned the discrepancy between that prior and measured elastic energy absorption. Three simulation-informed and three uninformative-prior campaigns each received a 30-experiment budget, with predicted optima tested using repeated fabricated samples~\cite{asd0615}. The informed method performed better early, but the approaches were not statistically distinguishable after 30 experiments. The increment was therefore faster acquisition of high-performing designs, rather than an established difference in final attainable performance. For toughness, the same straightforward transfer of the elastic prediction did not accelerate simulated search; useful information instead entered through a yield-force constraint and a transformation of the experimental response~\cite{asd0615}. Physics contributes when its representation matches the decision, even when it does not predict the target property directly.

Task changes expose a separate limitation in instrument-workflow comparison. A microscopy study evaluated 105 agent configurations over 53 tests. A random-forest surrogate predicting configuration performance had ROC-AUC 0.78 when holding out configurations, but 0.55 when entire test identities were held out~\cite{asd1009}. Selected agents nevertheless completed additional physical microscopy workflows. The failure concerns generalization of configuration-performance predictions to new tasks, not universal inability to operate the instrument; ROC-AUC is not the agent's task-success rate. A configuration effective for known workflows need not remain preferable under a new scientific task.

\subsection{Aerospace, Marine, and Transportation Sciences}

Mission systems provide longstanding examples of model-based planning under explicit temporal and resource constraints. The Autonomous Sciencecraft Experiment combines science analysis, planning and scheduling, and robust execution to adapt EO-1 observations~\cite{asd0645}. Its software was already commanding the spacecraft in orbit, although the conference demonstration also included a simulated science-response scenario. The distinction matters when attributing evidence to operational or simulated behavior.

Related rover studies describe OASIS planning, scheduling, execution, and plan modification in response to opportunities and execution conditions~\cite{asd0712,asd0733}. Evaluation includes simulation and terrestrial prototype hardware in the JPL Mars Yard; the 2005 implementation also describes components communicating with the rover from offboard workstations. These reports concern a related system lineage, rather than independent demonstrations of unrelated architectures.

Operators or onboard analysis provide goals, while planners and execution monitors allocate activities within permitted actions. Replanning establishes responsiveness under those models; autonomous question formation and additional scientific return require separate evidence. The selected spacecraft and rover studies support mission-science comparisons without implying equivalent evidence for every marine and transportation subfield.

\subsection{Social, Behavioral, Economic, and Knowledge-System Sciences}

Scientific communities and institutions can themselves be research objects. SciSciGPT uses language-based tool coordination to retrieve and analyze scholarly records in collaboration with researchers~\cite{asd0624}. Its construction of an inter-university collaboration network through database queries and iterative visualization illustrates empirical knowledge-production analysis. Here, the scientific evidence depends on data provenance and the adequacy of the analytical procedure.

ResearchTown models reading, idea generation, writing, and reviewing within a research community, evaluating reconstruction of masked papers and reviews against real reference outputs using text-similarity metrics~\cite{asd0652}. AgentReview uses interacting language-model agents and prompt-defined reviewer and chair characteristics to study peer-review dynamics; the reported simulated outcomes were not directly compared with actual peer-review decisions~\cite{asd0786}. Human specification of assumptions is therefore consequential for interpreting the generated interactions.

Scholarly-data analysis, reconstruction similarity, and exploration of behavioral assumptions address different propositions. Multiple interacting agents describe an organizational design; they do not independently validate the modeled institution. External empirical evidence is required to connect simulated behavior to human mechanisms. These knowledge-system examples form a specific part of the broader social, behavioral, and economic module.
